\begin{helper}[The sign-free core of the weighted Caccioppoli estimate at a level]
Let $n\ge2$, let $p>1$ and let $r>0$.  Let $\varepsilon,\gamma_0,\gamma_1\in\R$,
let $\Omega\subseteq\R^n$, let $A\colon\R^n\to\R^n$, let $u\colon\R^n\to\R$, let
$g\colon\R^n\to\R^n$, let $Dg$ assign to each index $i$ a field
$Dg^i\colon\R^n\to\R^n$, let $\zeta\colon\R^n\to\R$ and let $k$ be an index in
$\{1,\dots,n\}$.  Assume
\begin{enumerate}
\item[(H1)] $0<\varepsilon$, $0<\gamma_0$, $\gamma_1\le r\gamma_0$ and
  $\gamma_0\le\gamma_1$;
\item[(H2)] $\Omega$ is open and nonempty;
\item[(H3)] the field $(x,z,h)\mapsto A(h)$ satisfies (Q1) of the thesis and
  (Q2) of the thesis on $\Omega$ with the parameters
  $p,\varepsilon,\gamma_0,\gamma_1$ (\noderef{EllipticStructureQ12});
\item[(H4)] the cited regularity result (C1) holds for $A$ on $\Omega$ with
  those parameters (\noderef{CitedRegularityEpsPos});
\item[(H5)] $u$ is a weak solution of $\operatorname{div}A(\nabla u)=0$ in
  $\Omega$ with weak gradient $g$ (\noderef{IsWeakSolutionOn});
\item[(H6)] $g^i\in W^{1,2}_{\mathrm{loc}}(\Omega)$ with weak gradient $Dg^i$,
  for every index $i$ (\noderef{MemW1pLocOn});
\item[(H7)] $\zeta$ is $C^\infty$ with compact support and
  $\operatorname{supp}\zeta\subseteq\Omega$, where $\operatorname{supp}$ denotes
  the closed support.
\end{enumerate}
Let $W\colon\R^n\to\R$ satisfy
\begin{enumerate}
\item[(H8)] $W=\bigl(\varepsilon+|g|^2\bigr)^{\frac p2-1}$ pointwise, the power
  being the real power of the structure conditions, written exactly as
  (Q1) of the thesis--(Q2) of the thesis print it.  ($W$ enters as a function argument
  together with this defining equation, so that a consumer may pass its own
  copy.)
\end{enumerate}
Finally let $T\subseteq\R^n$, let $\sigma\in\R$, let $v\colon\R^n\to\R$ and let
$Z\colon\R^n\to\R^n$ satisfy
\begin{enumerate}
\item[(H9)] $\sigma=1$ or $\sigma=-1$;
\item[(H10)] $T$ is measurable and $T\subseteq\Omega$;
\item[(H11)] $v\ge0$ everywhere;
\item[(H12)] for almost every $x\in\Omega$: if $x\in T$ then
  $Z(x)=\sigma\,Dg^k(x)$, and if $x\notin T$ then $Z(x)=0$ and $v(x)=0$;
\item[(H13)] $v\in W^{1,2}_{\mathrm{loc}}(\Omega)$ with weak gradient $Z$
  (\noderef{MemW1pLocOn}).
\end{enumerate}
Then
\[
\int_T W\,|Dg^k|^2\zeta^2\,dx
\ \le\ 4n^2r^2\int_\Omega W\,v^2\sum_{i=1}^n(D_i\zeta)^2\,dx .
\]
\end{helper}
\begin{proof}
Write $K=\operatorname{supp}\zeta$, which is compact by (H7), and
\[
X=\int_TW|Dg^k|^2\zeta^2\,dx,
\qquad
Y=\int_\Omega W\,v^2\sum_i(D_i\zeta)^2\,dx ,
\]
both integrands being nonnegative, so that $X,Y\in[0,\infty]$; finiteness is
established in Step 3.  From (H1), $\gamma_1>0$, and $W>0$ everywhere, the base
$\varepsilon+|g|^2$ being at least $\varepsilon>0$.  Put
$M_{ij}=D_{h_j}A^i(g(\cdot))$ for the coefficient matrix.

\emph{Step 0: the structure conditions at every point, and $A\in C^1$.}  Fix a
point $x_\ast\in\Omega$, available by (H2), and read clause (4) of (H3) --- for
every $x\in\Omega$ and every $z$ the map $h\mapsto A(x,z,h)$ is of class $C^1$
on $D(\varepsilon)=\{h:\varepsilon+|h|^2\ne0\}$ --- at $x_\ast$ and $z=0$.
Since $\varepsilon>0$ \emph{strictly} by (H1), $D(\varepsilon)=\R^n$, so $A$ is
of class $C^1$ on all of $\R^n$, the symbols $D_{h_j}A^i$ are genuine partial
derivatives and $D_hA$ is continuous.  \emph{The strictness of
$\varepsilon>0$ is what makes this citation available}: at $\varepsilon=0$
clause (4) gives nothing at $h=0$, and for $1<p<2$ the admissible field
$A(h)=|h|^{p-2}h$ is not differentiable there.  Reading (Q1) of the thesis and
(Q2) of the thesis of (H3) the same way, and expanding
$(D_hA(g)\xi)_i=\sum_jM_{ij}\xi_j$ by linearity of $D_hA(g)$ in its direction
argument, we obtain for almost every $x\in\Omega$
\[
\gamma_0W|\xi|^2\le\sum_{i,j}M_{ij}\xi_j\xi_i\ \ (\xi\in\R^n),
\qquad
|M_{ij}|\le\gamma_1W\ \ (i,j) .
\tag{1}
\]

\emph{Step 1: the test pair.}  Put
\[
\varphi=v\,\zeta^2,
\qquad
g_\varphi=\zeta^2Z+2v\,\zeta\,\nabla\zeta .
\]
Since $\zeta$ is $C^\infty$ with compact support in $\Omega$, so is $\zeta^2$,
with $\nabla(\zeta^2)=2\zeta\nabla\zeta$ and
$\operatorname{supp}(\zeta^2)\subseteq K$.  By (H13) and the product rule for a
smooth compactly supported factor (\noderef{WeakGradientSmoothMul}, clause (2),
at $q=2$ and $\Omega$), $\varphi\in W^{1,2}(\Omega)$ with weak gradient
$g_\varphi$ (\noderef{MemW1pOn}), $\varphi$ has compact support and
$\operatorname{supp}\varphi\subseteq K\subseteq\Omega$.

\emph{Step 2: substitution into the differentiated equation.}  Apply
\noderef{DifferentiatedEquationWeak} with the present
$n,p,\varepsilon,\gamma_0,\gamma_1,\Omega,A,u,g,Dg$, the index $k$ and the pair
$(\varphi,g_\varphi)$ of Step 1.  Its hypotheses are $2\le n$, $1<p$,
$0<\varepsilon$, $0<\gamma_0$, $\gamma_0\le\gamma_1$ --- all of (H1) --- the
openness of $\Omega$, (H3), (H4), (H5) and (H6), together with the admissibility
of $(\varphi,g_\varphi)$, which is Step 1.  Its conclusion is
\[
\int_\Omega\sum_{i,j}M_{ij}\,(Dg^k)_j\,(g_\varphi)_i\,dx=0 .
\tag{2}
\]

\emph{Step 3: the two terms of (2), their measurability and their finiteness.}
Write $T_1$ and $T_2$ for the integrals obtained from (2) by replacing
$g_\varphi$ by its two summands $\zeta^2Z$ and $2v\zeta\nabla\zeta$.  Both
integrands vanish off $K$, because $\zeta$ does.  Measurability first: by (H4)
there is a field $g'$ continuous on $\Omega$ with $g=g'$ almost everywhere
there, so each $M_{ij}$ agrees almost everywhere on $\Omega$ with the
continuous $D_{h_j}A^i(g'(\cdot))$ --- using Step 0's continuity of $D_hA$ ---
and is therefore almost everywhere strongly measurable; the same two sentences
apply to $W$, whose $g'$-version $(\varepsilon+|g'|^2)^{p/2-1}$ is continuous on
$\Omega$, the base being continuous and bounded below by $\varepsilon>0$ and
$t\mapsto t^{p/2-1}$ continuous on $(0,\infty)$.  The field $Dg^k$ is almost
everywhere strongly measurable on $\Omega$, being a weak gradient (H6); $v$ and
$Z$ are, by (H13); and $\zeta$, $\nabla\zeta$ are continuous.  Every integrand
below is a finite sum of products of these, hence almost everywhere strongly
measurable on $\Omega$.

For the domination, by (1) and $|Z|\le|Dg^k|$ almost everywhere --- which
follows from (H12) and $|\sigma|=1$, $Z$ being $\sigma Dg^k$ on $T$ and $0$ off
it --- together with $\sum_j|a_j|\le\sqrt n\,|a|$ (Cauchy--Schwarz, Preamble)
applied twice and the elementary $2st\le s^2+t^2$,
\[
\Bigl|\zeta^2\sum_{i,j}M_{ij}(Dg^k)_j Z_i\Bigr|
\le n\gamma_1M_W\bigl(\sup\zeta^2\bigr)|Dg^k|^2 ,
\]
\[
\Bigl|2v\,\zeta\sum_{i,j}M_{ij}(Dg^k)_jD_i\zeta\Bigr|
\le n\gamma_1M_W\bigl(\sup2|\zeta||\nabla\zeta|\bigr)
\bigl(|Dg^k|^2+v^2\bigr) ,
\]
where $M_W$ is a bound for $W$ on the compact $K$, available because $W$ agrees
almost everywhere on $\Omega$ with a function continuous on $\Omega$ and
$K\subseteq\Omega$ is compact.  The functions $|Dg^k|^2$ and $v^2$ are
integrable on $K$ by (H6) and (H13) together with H\"older on a set of finite
measure (Preamble), and $\zeta$, $\nabla\zeta$ are continuous with compact
support, hence bounded.  So $T_1$ and $T_2$ are finite by domination
(Preamble), (2) splits as $T_1+T_2=0$, i.e.\ $T_1=-T_2$, and the same bounds
give $X<\infty$ and $Y<\infty$.

\emph{Step 4: the lower bound for $\sigma T_1$.}  By (H12) the integrand of
$T_1$ vanishes almost everywhere off $T$, where $Z=0$, and equals
$\sigma\,\zeta^2\sum_{i,j}M_{ij}(Dg^k)_j(Dg^k)_i$ almost everywhere on $T$,
where $Z=\sigma Dg^k$.  Multiplying by $\sigma$ and using $\sigma^2=1$, which
holds by (H9), the first inequality of (1) at $\xi=Dg^k(x)$ gives
\[
\sigma T_1
=\int_T\zeta^2\sum_{i,j}M_{ij}(Dg^k)_j(Dg^k)_i\,dx
\ \ge\ \gamma_0\int_TW|Dg^k|^2\zeta^2\,dx=\gamma_0X ,
\]
monotonicity of the integral (Preamble) applying because both integrands are
integrable by Step 3.

\emph{Step 5: the upper bound for $|T_2|$.}  The integrand of $T_2$ vanishes
almost everywhere off $T$, because $v$ does by (H12).  By the norm inequality
for the integral and the second inequality of (1) with two Cauchy--Schwarz
steps in the indices,
\[
|T_2|\le2n\gamma_1\int_TW\,v\,|\zeta|\,|Dg^k|\,|\nabla\zeta|\,dx .
\]
Apply the Cauchy--Schwarz inequality (Preamble) on $T$ with respect to the
measure $W\,dx$, to the two factors $|\zeta||Dg^k|$ and $v|\nabla\zeta|$:
\[
|T_2|\ \le\ 2n\gamma_1
\Bigl(\int_TW|Dg^k|^2\zeta^2dx\Bigr)^{1/2}
\Bigl(\int_TW\,v^2|\nabla\zeta|^2dx\Bigr)^{1/2}
\ \le\ 2n\gamma_1X^{1/2}Y^{1/2},
\]
the last step because the first integral is $X$ and the second is at most $Y$,
the integrand being nonnegative, $T\subseteq\Omega$ by (H10) and
$|\nabla\zeta|^2=\sum_i(D_i\zeta)^2$.

\emph{Step 6: conclusion.}  Combining Steps 4 and 5 with $T_1=-T_2$, and noting
$\sigma T_1=|\sigma T_1|=|T_2|$ because $\sigma T_1\ge\gamma_0X\ge0$ and
$|\sigma|=1$,
\[
\gamma_0X\ \le\ 2n\gamma_1X^{1/2}Y^{1/2} .
\tag{3}
\]
If $X=0$ the assertion holds outright, its right-hand side being nonnegative.
If $X>0$ then, $X$ being finite by Step 3, divide (3) by $X^{1/2}>0$ to get
$\gamma_0X^{1/2}\le2n\gamma_1Y^{1/2}$, square, and divide by $\gamma_0^2>0$:
\[
X\ \le\ 4n^2\Bigl(\frac{\gamma_1}{\gamma_0}\Bigr)^2Y\ \le\ 4n^2r^2Y ,
\]
the second inequality because $0<\gamma_1/\gamma_0\le r$ by (H1) --- which gives
both $\gamma_0\le\gamma_1$ and $\gamma_1\le r\gamma_0$ --- so that
$(\gamma_1/\gamma_0)^2\le r^2$, and because $Y\ge0$.  That is the assertion.
\end{proof}
